\section{Distribución hipergeométrica}
\subsection{Definición}
La distribución de probabilidad de la variable aleatoria hipergeométrica $x$, el número de éxitos en una muestra aleatoria de tamaño $n$, seleccionada de $N$ resultados de los cuales $K$ son éxitos y $N-K$ son fracasos es:
\begin{equation}
    P(x;N,n,K) = \frac{C\left(\begin{matrix}
            K \\
            x
        \end{matrix}\right)C\left(\begin{matrix}
            N-K \\
            n-x
        \end{matrix}\right)}{C\left(\begin{matrix}
            N \\
            n
        \end{matrix}\right)} \qquad x \in \mathbb{N}
\end{equation}
\subsection{Propiedades}
\begin{itemize}
    \item Una muestra aleatoria de tamaño n se selecciona de N resultados.
    \item K de los N resultados se pueden clasificar como éxitos y N-K como fracasos.
\end{itemize}
De la definición se obtiene los siguientes resultados:
\begin{itemize}
    \item $\mu=\frac{nK}{N}$
    \item $\sigma^2= \left(\frac{N-K}{N-1}\right)\left(\frac{K}{N}\right)\left(1-\frac{K}{N}\right)n$
\end{itemize}
\begin{ejemplo}
    Un cargamento de 20 grabadoras magnetofónicas contiene 5 defectuosas. Si 10 de ellas son aleatoriamente escogidas para revisión: ¿Cuál es la probabilidad de que dos estén defectuosas?\\
    x=2, n=10, N=20, K=5
    \begin{align*}
        P(2;20,10,5) & =\frac{C\left(\begin{matrix}
                K \\
                x
            \end{matrix}\right)C\left(\begin{matrix}
                N-K \\
                n-x
            \end{matrix}\right)}{C\left(\begin{matrix}
                N \\
                n
            \end{matrix}\right)}     \\
                     & = \frac{C\left(\begin{matrix}
                5 \\
                2
            \end{matrix}\right)C\left(\begin{matrix}
                20-5 \\
                10-2
            \end{matrix}\right)}{C\left(\begin{matrix}
                20 \\
                10
            \end{matrix}\right)}  \\
                     & = \frac{C\left(\begin{matrix}
                5 \\
                2
            \end{matrix}\right)C\left(\begin{matrix}
                15 \\
                8
            \end{matrix}\right)}{C\left(\begin{matrix}
                20 \\
                10
            \end{matrix}\right)} \\
                     & = \frac{\frac{5!}{2!3!}\frac{15!}{8!7!}}{\frac{20!}{10!10!}}                                                                        \\
                     & = \frac{10(6435)}{184756}                                                                                                           \\
                     & = 0.3483
    \end{align*}
\end{ejemplo}